\documentclass[11pt]{article}
\usepackage[a4paper,margin=28mm]{geometry}
\usepackage{amsmath,amssymb,amsthm,xcolor,booktabs,array}
\newcommand{\PROVED}{\textcolor{green!40!black}{\textbf{PROVED}}}
\newcommand{\OPEN}{\textcolor{red!70!black}{\textbf{OPEN}}}
\newcommand{\AUDIT}{\textcolor{orange!70!black}{\textbf{AUDIT}}}
\newtheorem{theorem}{Theorem}
\newtheorem{proposition}{Proposition}
\newtheorem{lemma}{Lemma}
\title{The Witt--Tate adjoint and the odd Poisson quotient (v96)}
\date{13 August 2026}

\begin{document}
\maketitle

\section{The exact Tate renormalization}

Let $W_{\rm alg}=\bigoplus_{r\geq1}\mathbb C\phi_r$ be Haran's
cyclotomic Witt orbit.  Haran's state gives
\[
 \langle\phi_r,\phi_s\rangle_H=\delta_{rs}\varphi(r),
 \qquad V_n\phi_r=\phi_{nr},
\]
where $\varphi$ is Euler's totient.  Thus the statement in v92 that this
basis is orthonormal is withdrawn and classified \AUDIT.

Normalize first by Haran's state and then by the self-dual metric character
of row (b).  Equivalently put
\begin{equation}\label{eq:tatenorm}
 \langle\phi_r,\phi_s\rangle_{1/2}:=\frac{\delta_{rs}}{r}.
\end{equation}
Indeed, inside Haran's Hilbert space this is obtained from the diagonal
source-defined operator
\[
 C\phi_r=\frac1{\sqrt{r\varphi(r)}}\phi_r,
 \qquad
 \langle C\phi_r,C\phi_s\rangle_H=\frac{\delta_{rs}}r.
\]
The factor $\varphi(r)^{-1/2}$ normalizes the cyclotomic state and
$r^{-1/2}$ is precisely the half-Tate metric of row (b).

\begin{theorem}[Witt--Tate adjoint --- \PROVED]\label{thm:wittadjoint}
For the inner product \eqref{eq:tatenorm},
\begin{equation}\label{eq:Vadjoint}
 V_n^\dagger\phi_k=
 \begin{cases}
  n^{-1}\phi_{k/n},&n\mid k,\\
  0,&n\nmid k.
 \end{cases}
\end{equation}
Consequently $S_n:=\sqrt n\,V_n$ is an isometry.  Its minimal group dilation
on the completion of
$\bigoplus_{q\in\mathbb Q_+^\times}\mathbb C\phi_q$ with
$\|\phi_q\|^2=q^{-1}$ is unitary and satisfies
\[
 S_q^*=S_{q^{-1}}\qquad(q\in\mathbb Q_+^\times).
\]
\end{theorem}

\begin{proof}
For $k=nr$ one has
\[
 \langle V_n\phi_r,\phi_k\rangle_{1/2}=\frac1{nr}.
\]
If $V_n^\dagger\phi_{nr}=c\phi_r$, the other side of the adjoint identity is
$c/r$, hence $c=n^{-1}$.  If $n\nmid k$, orthogonality gives zero.  Also
\[
 \|S_n\phi_r\|_{1/2}^2
 =n\|\phi_{nr}\|_{1/2}^2=\frac1r.
\]
On positive rationals the same formula is onto and therefore unitary; the
adjoint is its inverse.
\end{proof}

\section{Transport to the Dirichlet and scaling realizations}

Give the algebraic Dirichlet coefficient space the norm
$\|\delta_r\|^2=r^{-1}$.  The row-(b) map $J(\phi_r)=\delta_r$ is then
unitary and intertwines $V_n$ with convolution by $\delta_n$.

Let $U_nh(x)=h(nx)$ on $L^2(\mathbb R_+,dx)$.  A change of variables gives
\begin{equation}\label{eq:scalingadjoint}
 U_n^*=n^{-1}U_{1/n}.
\end{equation}

\begin{theorem}[Algebraic Witt--Rosati transport --- \PROVED]
On the common algebraic scaling domain, the row-(b)/(c) realization carries
the Witt adjoint to the Tate adjoint:
\[
 J V_n^\dagger J^{-1}=n^{-1}U_{1/n}=U_n^*.
\]
Equivalently, for the group algebra generated by rational dilations, the
involution transported from the Tate-renormalized Witt state is the
$L^2(dx)$ adjoint.
\end{theorem}

\begin{proof}
Both sides send the vector labelled $nr$ to $n^{-1}$ times the vector
labelled $r$ and vanish on labels not divisible by $n$.  Formula
\eqref{eq:scalingadjoint} proves the analytic equality.
\end{proof}

This proves the missing adjoint comparison \emph{before} passage to Meyer's
odd quotient.  It uses no zero, no explicit-formula sign, and no choice of a
positive part.

\section{Descent to the odd quotient}

Put
\[
 \mathcal H_-^0=\mathcal H_-/Z\mathcal H_\cap.
\]
The subspace $Z\mathcal H_\cap$ is invariant under all rational dilations and
their algebraic Tate adjoints.  Hence the involution descends algebraically
to the quotient.  This statement alone is not yet a Hilbert adjoint theorem:
one needs a positive completion in which the subspace is closed (or,
equivalently, in which its closure does not enlarge the radical relevant to
the nuclear character).

\begin{proposition}[Exact descent criterion --- \PROVED]
Assume there is a positive Hilbert seminorm $\|\cdot\|_-$ on $\mathcal H_-$
such that:
\begin{enumerate}
\item its null space is zero and the rationally normalized scalings
      $\sqrt n\,U_n$ extend isometrically;
\item $Z\mathcal H_\cap$ is closed for this norm;
\item the Hilbert quotient action has the same nuclear character as the odd
      Fr\'echet quotient of row (c).
\end{enumerate}
Then for every test function $f$,
\[
 \pi_-(f^\vee)=\pi_-(f)^*,
\]
and for primitive $f$,
\[
 B_{\rm nuc}(f,f)
 =-\operatorname{Tr}\bigl(\pi_-(f)\pi_-(f)^*\bigr)\leq0.
\]
Thus row (d) and condition $G$ follow.
\end{proposition}

\begin{proof}
By Theorem~\ref{thm:wittadjoint}, the adjoint identity holds on a dense
algebraic domain.  Closedness and invariance make the quotient a reducing
Hilbert quotient, so the identity descends and extends by continuity to the
test algebra.  The even supertrace is
$|\widehat f(0)|^2+|\widehat f(1)|^2$ and vanishes on primitive $f$.
Character compatibility then turns the odd part of the row-(c) supertrace
into the displayed negative Hilbert--Schmidt square.
\end{proof}

\section{Why the ordinary $L^2(dx)$ quotient is not enough}

The ambient identity \eqref{eq:scalingadjoint} is exact, but merely completing
in $L^2(dx)$ cannot be used without checking the quotient.  The two Tate
moment conditions are not closed conditions in ordinary $L^2$, and the
Poisson range can acquire a larger closure.  If its closure is all of the
ambient space, the Hilbert quotient is zero and no longer has Meyer's nuclear
character.  Therefore ``take the $L^2$ quotient'' is \AUDIT unless clauses
(2)--(3) of the descent criterion are proved.

Likewise, the local Witt Hilbert completion in v92 used the wrong
orthonormality.  The corrected norm \eqref{eq:tatenorm} fixes the adjoint but
does not by itself make the global Zeta/Poisson range closed.

\section{Terminal ledger}

\begin{center}
\begin{tabular}{>{\raggedright\arraybackslash}p{.48\linewidth}
                >{\raggedright\arraybackslash}p{.18\linewidth}
                >{\raggedright\arraybackslash}p{.24\linewidth}}
\toprule
Statement & Status & Location \\
\midrule
Canonical Tate-renormalized Witt norm $\|\phi_r\|^2=1/r$ & \PROVED &
\eqref{eq:tatenorm}\\
Exact generator adjoint $V_n^\dagger=n^{-1}V_{1/n}$ & \PROVED &
Theorem~\ref{thm:wittadjoint}\\
Compatibility with the scaling adjoint $U_n^*=n^{-1}U_{1/n}$ & \PROVED &
Algebraic transport theorem\\
Algebraic descent of the involution to $\mathcal H_-^0$ & \PROVED &
Invariance of $Z\mathcal H_\cap$\\
Closed positive Hilbert descent with preservation of the nuclear character &
\OPEN & Clauses (2)--(3) of the descent criterion\\
Global row (d) and condition $G$ & \OPEN & Immediate after the preceding
descent theorem\\
\bottomrule
\end{tabular}
\end{center}

\end{document}
